\section*{NS94: tangent-circle inversions and what they can detect}
\paragraph{Source and scope.} Yufei Zhao's four-page 2011 handout
\emph{Power of a Point}, supplied by the user, defines power and radical
axes and proves the radical-axis theorem. It does not define exterior
similarity points or discuss zeta zeros. We derive the similarity formula
below. No instructions in the handout are treated as project instructions.
All controls below are exact geometric examples, not hypothetical zeros
asserted to exist for zeta. RH remains open.

\paragraph{A line under one inversion.}
With real centre $c$ and radius $r>0$, ordinary circle inversion is
$I_{c,r}(z)=c+r^2/(\bar z-c)$. The exceptional centre maps to infinity.
For a vertical line $x=a$, $a\ne c$, its image is the circle of centre
$m=c+r^2/[2(a-c)]$ and radius $R=r^2/[2|a-c|]$, with the limiting point
$c$ supplied only by the point at infinity of the original line.
More generally, for $z=x+iy\ne c$,
\[
 |I_{c,r}(z)-m|^2-R^2
 =\frac{r^4(a-x)}{(a-c)((x-c)^2+y^2)}.
\]
This identity follows by expanding the two squares; the worksheet checks it.

\paragraph{One inversion in each tangent circle.}
Take $a=1/2$, centres $a-p,a+q$ and radii $p,q>0$. This is the externally tangent interpretation, with one circle on each side. Both circles
are tangent to the critical line at $a$. The two image circles have centres
$a-p/2,a+q/2$ and radii $p/2,q/2$, respectively. They lie within the
original circles and touch them at $a$.
For distinct image radii the exterior homothety centre is
\[
 e=\frac{R_Lm_R-R_Rm_L}{R_L-R_R}=a+\frac{pq}{p-q}.
\]
This follows by solving $m_R=e+(R_R/R_L)(m_L-e)$; it contains no zeros.
For $p=q$ and distinct centres there is no finite exterior homothety centre;
it is the point at infinity along their line of centres. This already
happens with no marked points on either circle.
Shrinking both radii does not determine a unique limit: for $p=t,q=2t$,
$e=a-2t\to a$; for $p=t,q=t+t^2$, $e=a-1-t\to a-1$;
for $p=t,q=t+t^3$, $e=a-t^{-1}-t\to-\infty$.
All three paths use the same geometric construction, independently of zeros.
For internally tangent circles on the same side, centres $a-p,a-q$ and
radii $p,q$, the image centres are $a-p/2,a-q/2$ and the exterior
homothety centre is simply $a$ when $p\ne q$. The proposed count
dependence does not arise in that interpretation either.

\paragraph{Two successive inversions: the other interpretation.}
Write $u=z-a$. Applying the left inversion and then the right gives
\[
 I_R I_L(z)-a=\frac{u}{1+ku},\qquad k=p^{-1}+q^{-1}.
\]
Reversing the order gives $u/(1-ku)$. Each maps the critical line
onto a circle (with its missing limiting point): centres $a\pm1/(2k)$,
common radius $1/(2k)$. Thus doing both orders produces equal circles
whose exterior similarity point is again at infinity, for every $p,q>0$.
A single composition gives one image circle, not two.

\paragraph{What an off-line point does.}
Let $P_L,P_R$ denote powers relative to the two separately inverted image
circles. Then
\[
 P_L(I_Lz)=\frac{p^3(a-x)}{(x-a+p)^2+y^2},\qquad
 P_R(I_Rz)=\frac{q^3(x-a)}{(x-a-q)^2+y^2}.
\]
Off-line points therefore go inside one image circle and outside the other.
This detects location; it does not imply an imbalance in zero counts.
These are images of zero locations, not a claim that $\zeta(I(z))=0$
whenever $\zeta(z)=0$. Inversion supplies no new zeta functional equation.
The nontrivial zeta zeros are symmetric under $z\mapsto1-\bar z$,
with multiplicities. Reflection supplies an opposite-side partner at the
same height. Counts must use the same finite original height cutoff;
comparing different image cutoffs compares different sets, and comparing
two unqualified infinite counts is not a finite discrepancy test.

An exact control is
\[
 P(s)=((s-a-\delta)^2+t^2)((s-a+\delta)^2+t^2),\quad
 a=1/2,\quad\delta=1/4,\quad t=1.
\]
It is real and $P(1-s)=P(s)$, with four off-line roots
$a\pm\delta\pm it$. For $p=q=1/2$, each image disc contains two
of their transformed roots; each image boundary contains none. The
exterior similarity point is infinite, just as for equal circles marked
only with critical-line points. This refutes the proposed geometric
count-imbalance inference. The polynomial lacks zeta's arithmetic data;
it is not a counterexample to RH or a universal exclusion of geometry.

\paragraph{A usable location statistic, but an unproved target.}
For the same point, the product yields
\[
 -\frac{P_L(I_Lz)P_R(I_Rz)}{p^3q^3}
 =\frac{(x-a)^2}{[(x-a+p)^2+y^2][(x-a-q)^2+y^2]}\ge0.
\]
For the actual nontrivial zeros within a fixed finite height $H$, summing
this quantity with multiplicities gives a nonnegative number, zero exactly
when every included zero lies on the critical line. It avoids cancellation
between reflected off-line partners. Proving it zero for every height is
an RH-equivalent target, not a consequence of inversion or symmetry.
No arithmetic estimate for this statistic is supplied.

\paragraph{The proposed variant at 0 and 1.}
The tangency partners were not specified. Under the natural inward-facing
choice of centres $r,1-r$ and equal radius $r<1/2$ (tangent to the vertical
lines at 0 and 1), the critical-line image radii are
$r^2/[2(1/2-r)]\to\infty$ as $r\uparrow1/2$. At the endpoint the
critical line passes through both inversion centres and is preserved as a
line. Under the outward-facing choice $-r,1+r$, they tend instead to
$1/8$. These are explicitly stated interpretations, not an assertion about
an unspecified construction.

\paragraph{Relation to the earlier project.}
Clifford algebra in NS47 represents the existing signed form; it did not
supply its missing uniform lower bound. The complete first-zero enclosure
is a fixed-window result, not control of all zeta zeros. The current NB
work constructs approximations and seeks to prove their complete error
tends to zero. No Clifford machinery is needed for this inversion check.
